\documentclass[11pt,a4paper]{article}
\usepackage{avaliacao-poo}
\configurarAvaliacao{Checkpoint [número]}{Programação Orientada a Objetos}
\validarCheckpointDuasPaginas
\begin{document}
\cabecalhoAvaliacao
\begin{quadroRespostas}
\linhaRespostas{Questão 1}{a,b,c}
\linhaRespostas{Questão 2}{a,b,c}
\end{quadroRespostas}
\cabecalhoCenario{[Nome do cenário]}
[Contexto curto com \codigoInline{codigo.inline()}.]
\noindent\begin{minipage}[t]{.39\textwidth}
\begin{painelreferencia}{Classe.java}
\begin{lstlisting}[style=cenariocodigo]
public class Classe {
    int estado;
    void agir() { estado = estado + 1; }
}
\end{lstlisting}
\end{painelreferencia}
\end{minipage}\hfill
\begin{minipage}[t]{.59\textwidth}
\begin{painelreferencia}{Main.java}
\begin{lstlisting}[style=cenariocodigo]
public class Main {
    public static void main(String[] args) {
        Classe objeto = new Classe();
        objeto.agir();
        System.out.println(objeto.estado);
    }
}
\end{lstlisting}
\end{painelreferencia}
\end{minipage}
\inicioQuestoes
\questaoDireta{1}{[Enunciado objetivo e determinado.]}{[pontos]}
\newpage
\questaoDireta{2}{[Enunciado objetivo e determinado.]}{[pontos]}
\begin{painelEvolucao}{Trecho alterado}
\begin{lstlisting}[style=evolucaocodigo]
private int estado;
public int getEstado() { return estado; }
\end{lstlisting}
\end{painelEvolucao}
\end{document}
